Since $\operatorname{d\acute er}(G)$ is smooth over $S$ and has connected fibers, then,{\renewcommand{\thefootnote}{(64)}\nde{The original has been expanded in what follows. In particular, the conclusion (implicit in the original) that every subgroup of $G$, smooth with connected fibers and containing $\operatorname{d\acute er}(G)$, is closed in $G$, has been added.}} by Exp.~IV, 5.3.1 and 6.3.1, and Exp.~IVB 9.2, the map $H\mto f_0(H)$ establishes a bijective correspondence between group subschemes (resp. closed group subschemes) of $G$, containing $\operatorname{d\acute er}(G)$, smooth over $S$ and with connected fibers, and group subschemes (resp. closed group subschemes) of $\operatorname{corad}(G)$, smooth over $S$ and with connected fibers.
% typo?: kept as printed: Exp. IVB 9.2.

Now, if $H'$ is a group subscheme of $\operatorname{corad}(G)$, smooth over $S$ with connected fibers, then $H'$ is of finite presentation over $S$ (Exp.~VIB 5.5) and its fibers are tori (since those of $\operatorname{corad}(G)$ are), hence by Exp.~X 8.2, $H'$ is a subtorus of $\operatorname{corad}(G)$, hence is closed in $\operatorname{corad}(G)$ (Exp.~IX 2.6).

Consequently, every subgroup of $G$, smooth with connected fibers and containing $\operatorname{d\acute er}(G)$, is closed in $G$.{\renewcommand{\thefootnote}{(64)}\footnotemark[65]}
\numpara{6.3.2}\label{XXII.6.3.2}
If $H$ is a closed group subscheme of $G$, smooth over $S$, with connected fibers and invariant in $G$, then $H$ is reductive. If moreover $H\supset\operatorname{d\acute er}(G)$, then $\operatorname{d\acute er}(H)=\operatorname{d\acute er}(G)$ and $f_0(H)$ is identified with $\operatorname{corad}(H)$. One has therefore proved the
\begin{proposition}{6.3.3}\label{XXII.6.3.3}
Let $G$ be a reductive $S$-group. Every group subscheme $H$ of $G$, invariant in $G$, with commutative quotient (i.e. containing $\operatorname{d\acute er}(G)$), smooth over $S$, with connected fibers,{\renewcommand{\thefootnote}{(65)}\nde{The hypothesis that $H$ is retrocompact in $G$ has been deleted, since it holds automatically: by VIB 5.5, $G$ and $H$ are separated and quasi-compact over $S$, hence $H\hookrightarrow G$ is quasi-compact by EGA IV$_1$, 1.1.2 (v).}} is closed and reductive. One has $\operatorname{d\acute er}(H)=\operatorname{d\acute er}(G)$ and $f_0(H)$ is identified with $\operatorname{corad}(H)$; one has
\[
G/H\simeq\operatorname{corad}(G)/\operatorname{corad}(H),\qquad H=(H\cap\operatorname{rad}(G))\cdot\operatorname{d\acute er}(G).
\]
Moreover, $H\mto f_0(H)$ defines a bijection between the set of subgroups $H$ of $G$ having the preceding properties and the set of subtori of $\operatorname{corad}(H)$.
% typo?: kept as printed: subtori of corad(H), rather than corad(G).
\end{proposition}
By a new application of Noether's theorem (Exp.~IV, 5.3.1 and 6.3.1), one deduces the
\begin{proposition}{6.3.4}\label{XXII.6.3.4}
Let $S$ be a scheme, $G$ a reductive $S$-group, $T$ a maximal torus of $G$. For every subgroup $H$ of $G$ as above, $T\cap H$ is a maximal torus of $G$ and one has
\[
G/H\simeq T/T\cap H,\qquad H=(T\cap H)\cdot\operatorname{d\acute er}(G).
\]
Moreover, $H\mto T\cap H$ is a bijection between the set of subgroups $H$ of $G$ as above and the set of subtori of $T$ containing $T\cap\operatorname{d\acute er}(G)$.
% typo?: kept as printed: T intersect H is a maximal torus of G.
\end{proposition}

\begin{thebibliography}{DG70}
\bibitem[BLie]{XXII.BLie} N.~Bourbaki, \textit{Groupes et alg\`ebres de Lie}, Ch.~IV--VI, Hermann, 1968.
\bibitem[Ch05]{XXII.Ch05} C.~Chevalley, \textit{Classification des groupes alg\'ebriques semi-simples} (with the collaboration of P.~Cartier, A.~Grothendieck, M.~Lazard), Collected Works, vol.~3, Springer, 2005.
\bibitem[DG70]{XXII.DG70} M.~Demazure, P.~Gabriel, \textit{Groupes alg\'ebriques}, Masson \& North-Holland, 1970.
\bibitem[Gi71]{XXII.Gi71} J.~Giraud, \textit{Cohomologie non ab\'elienne}, Springer-Verlag, 1971.
\bibitem[Se64]{XXII.Se64} J.-P.~Serre, \textit{Cohomologie galoisienne}, Springer-Verlag, 1964; 5th ed.~1994.
\end{thebibliography}
